\section{Initial spectra and rigidity of rank-optimal receivers}
\label{sec:spectralrigidity93}
The dimension spectrum determines an optimal value.  We now determine
which initial spectra can attain it, and how a small loss in value
restricts the acquisition marginals.  Throughout this section the basis
index is drawn once and retained for both calls, the hypothesis priors
are equal, and $0<t\le1$.  The family satisfies
\eqref{eq:basismoments90}.  Put
\[
 r=\min\{k,\ell\},\qquad
 p_r=\frac12+\frac{t^2(dr-1)}{2dr(d+1)},\qquad
 \alpha=\frac{t^2}{2d^2(d+1)}.
\]
All ranks and norms refer to the fixed input copies.  Trace norms are
unhalved.  An initial Gram matrix $\rho$ means a pure first acquisition
$|C_0\rangle$ with $C_0^*C_0=\rho$, in the convention of
Lemma~\ref{lem:resetnormal91}; its input marginal is $\rho^{\mathsf T}$.
The receiver cut is exactly Definition~\ref{def:rankclass92}.

\subsection{The complete set of optimal initial spectra}
\begin{theorem}[Optimal initial Gram matrices]\label{thm:initialspectra93}
A fixed initial Gram matrix $\rho\in\mathcal D_d$ admits a strategy
in $\mathfrak R_{k,\ell}$ attaining $p_r$ if and only if
\begin{equation}\label{eq:spectralcap93}
                  \rho\preceq I/r.
\end{equation}
Every such $\rho$ has a decomposition
\begin{equation}\label{eq:projectionframe93}
 \rho=\sum_{h=1}^{s}w_hP_h/r,\qquad s\le d,\quad
 w_h>0,\quad\sum_hw_h=1,
\end{equation}
where the $P_h$ are rank-$r$ projections commuting with $\rho$.
This decomposition gives an attaining strategy with the same receiver
instrument after every first device label and with both retained
quantum registers of dimension $r$.

Except when $(d,r)=(2,1)$, every attaining strategy in the
dimension-preserving refined normal form has, on each nonzero branch,
\begin{equation}\label{eq:optimalgrams93}
       A_{yh}=w_{yh}P_{yh}/r,\qquad b_{yh}=P_{yh}/r,
       \qquad \operatorname{rank}P_{yh}=r.
\end{equation}
Together with the common Gram sums and maximizing leaf decisions,
these conditions are also sufficient.  They specify Gram data, not
unique dilations or unique final measurements.
\end{theorem}
\begin{proof}
For a normalized branch write $a=A/\tr A$ and $b=b_{yh}$.
The equal-prior upper in Theorem~\ref{thm:rankspectrum92}, before
applying its final inequality, is
\begin{equation}\label{eq:unsaturatedrankupper93}
 v\le\frac12+\alpha\sum_{y,h}w_{yh}
       \| (\sqrt{a_{yh}}\otimes\sqrt{b_{yh}})(I-dS)
                    (\sqrt{a_{yh}}\otimes\sqrt{b_{yh}})\|_1.
\end{equation}
The smaller Gram rank is at most $r$, and $\sum_{y,h}w_{yh}=d$.
All branch deficits from the bound $d-1/r$ are nonnegative.
Attainment therefore forces equality on every branch of positive
weight.  When $(d,r)\ne(2,1)$, Lemma~\ref{lem:rankswap92} gives
\eqref{eq:optimalgrams93}.  For each fixed $y$, summing yields
$\rho=\sum_hw_{yh}P_{yh}/r\preceq I/r$, since
$\sum_hw_{yh}=1$.  In the exceptional case $r=1$, the claimed spectral
condition holds for every density matrix anyway.  The fixed-$\rho$
compact-envelope argument of Theorem~\ref{thm:rankroof92} gives
attainment of its constrained value, so the same conclusion applies
when an optimum was initially defined through tester closure.

We give a finite construction for the converse.  Diagonalize $\rho$
with eigenvalues $\lambda_i$ and set $x_i=r\lambda_i$ and
$s_i=\sum_{j\le i}x_j$, $s_0=0$.  Then $0\le x_i\le1$ and
$s_d=r$.  For almost every $u\in[0,1)$, mark the $r$ points
$u,u+1,\ldots,u+r-1$ in $[0,r)$.  Select coordinate $i$ if its
half-open interval $[s_{i-1},s_i)$ contains a marked point.
Each interval has length at most one, so exactly $r$ distinct
coordinates are selected.  The probability of selecting coordinate
$i$ is $x_i$.  The selected subset is constant between successive
fractional parts of the $s_i$.  There are at most $d$ such intervals
of positive length, because $s_0$ and $s_d$ have the same fractional
part.  Their lengths are the weights $w_h$, and their selected
coordinate projections are $P_h$.  Taking diagonal entries proves
\eqref{eq:projectionframe93}.  Endpoints have zero measure and may
be assigned either adjacent choice.  Rational eigenvalues give rational
weights by this construction; no rational eigenbasis is presumed.

Choose isometries $J_h:\C^r\to\C^d$ with $J_hJ_h^*=P_h$ and put
\[
 C_{yh}=\sqrt{w_h/r}\,J_h^*,\qquad
 D_{yh}=J_h^*/\sqrt r,\qquad V_{yh}=C_{yh}C_0^+.
\]
On the initial Schmidt support, the common Gram sum gives
$\sum_hV_{yh}^*V_{yh}=I$ and $V_{yh}C_0=C_{yh}$.
Thus $V_{yh}$ is a recorded instrument, and $|D_{yh}\rangle$ is an
independently prepared normalized fresh vector.  The same maps may be
used for every $y$.  The receiver centered swap is
$w_h(I_{r^2}-dS_r)/r^2$.  Declare the scalar hypothesis on the
antisymmetric receiver projection for equal labels and on the symmetric
projection for unequal labels.  Complementary effects complete the
measurement.  These decisions attain each positive-part bound, proving
sufficiency, including $r=1$.  Applying the same construction without
a prescribed commuting eigenbasis proves sufficiency of
\eqref{eq:optimalgrams93} with the common Gram sums.  Countable
normal forms are covered by the positive trace tails of
Lemma~\ref{lem:countable91}.
\end{proof}
In particular the initial state need not be maximally mixed when
$r<d$.  Its largest Schmidt weight must be at most $1/r$.
When $r=d$, the condition reduces to $\rho=I/d$.
The finite interval construction concerns a known initial spectrum,
not construction of the second-moment family itself.

\subsection{Counting the first acquisition as well}
Let $\mathfrak R_{q;k,\ell}$ impose the additional condition that the
entire reference accompanying the first input has dimension at most
$q$, conditional on the public preparation record.  Mixed states on
$I_1\otimes\C^q$ are allowed.  The old receiver after its recorded
instrument has dimension at most $k$ and the independent second
reference at most $\ell$, as before.  Workspaces discarded before a
specified cut do not enlarge a counted register.  Classical records
and the final joint receiver measurement have the same conventions as
Definition~\ref{def:rankclass92}.
\begin{corollary}[Three specified dimension cuts]\label{cor:threecuts93}
For $1\le q,k,\ell\le d$, put $s=\min\{q,k,\ell\}$.
On the same finite family the complete classes
$\mathfrak R_{q;k,\ell}$ have exact optimal values
\begin{equation}\label{eq:threecuts93}
 \begin{split}
 P_{q;k,\ell}^{\rm eq}
      &=\frac12+\frac{t^2(ds-1)}{2ds(d+1)},\\
 P_{q;k,\ell}^{\rm b}
      &=\frac{d+1}{2(d+1)-t^2}
          +\frac{t^2(s-1)}{2s[2(d+1)-t^2]}.
 \end{split}
\end{equation}
The second line uses exactly the prior of
Theorem~\ref{thm:operationalhierarchy90}.  Both optima are attained by
two independently prepared maximally entangled vectors on any fixed
rank-$s$ input subspace, with no intervening receiver instrument and
no feedback.  No first-call reference of dimension $d$ is required
unless $s=d$.
\end{corollary}
\begin{proof}
Resolve a mixed first preparation into a finite pure-state mixture in
$I_1\otimes\C^q$ and record its mixture index classically.  This
recovers its old score if the index is ignored, without adding a quantum
purification.  For each component $\operatorname{rank}C_0\le q$.
The Kraus refinement of Lemma~\ref{lem:ranknormal92} then gives
$\operatorname{rank}A_{yh}\le\min\{q,k\}$, while
$\operatorname{rank}b_{yh}\le\ell$.  The two rank-sensitive swap
bounds apply with $s$.  The trace sums, constant-decision baselines
and positive-part optimizations in Theorem~\ref{thm:rankspectrum92}
give both upper bounds.  Averaging over initial mixtures and public
seeds, and then taking norm limits, preserves these bounds.

For attainment fix an isometry $J:\C^s\to\C^d$ and use
$C_0=D=J^*/\sqrt s$.  Keep the first receiver unchanged and prepare
the second vector freshly.  For every first label the old Gram and
the fresh Gram are both $JJ^*/s$.  The two final spectral readouts
used in Theorem~\ref{thm:rankspectrum92} attain the bounds.  All
three counted registers have dimension $s$.  This proves the assertion
also at $s=1$, where the antisymmetric receiver projection is zero.
\end{proof}
The three bounds apply at specified cuts.  They are not a cap on the
total workspace of an arbitrary implementation or a removal of the
reset architecture.  In particular, saturation remains the reset
value, not the unrestricted coherent-acquisition value.

\subsection{A stable common-projection theorem}
For density matrices $a,b$ with smaller rank at most $r$, set
\[
 \Delta_{d,r}(a,b)=d-r^{-1}
  -\|(\sqrt a\otimes\sqrt b)(I-dS)(\sqrt a\otimes\sqrt b)\|_1.
\]
When $(d,r)\ne(2,1)$ define
\begin{equation}\label{eq:rankrigidityconstant93}
 c_{d,r}=d-1-r^{-1}>0,\qquad
 K_{d,r}=r+\frac{(3+\sqrt2)^2}{c_{d,r}}.
\end{equation}
\begin{lemma}[Quantitative common support]\label{lem:projectionstability93}
For $d\ge2$, $1\le r\le d$ and $(d,r)\ne(2,1)$, there exists
a rank-$r$ projection $P$ such that
\begin{equation}\label{eq:projectionstability93}
 \max\{\|a-P/r\|_1^2,\|b-P/r\|_1^2\}
       \le K_{d,r}\Delta_{d,r}(a,b).
\end{equation}
Only the smaller Gram rank need be bounded by $r$; neither
invertibility nor equality of the two supports is assumed.
\end{lemma}
\begin{proof}
Put $f=\|\sqrt a\sqrt b\|_1$, $\eta=1-f^2$ and
$e=\tr(ab)-f^2/r$.  The overlap rank bound gives $e\ge0$, and
$f\le1$ gives $\eta\ge0$.  The proof of
Lemma~\ref{lem:rankswap92} gives
\begin{equation}\label{eq:rankdefects93}
              \Delta_{d,r}(a,b)\ge c_{d,r}\eta+e.
\end{equation}
Extend a right polar factor to a unitary $U$ so that
$Z=\sqrt a\sqrt bU=\sqrt{\sqrt a b\sqrt a}\succeq0$.
The Hilbert--Schmidt identity and the Schatten product inequality give
\[
 \|\sqrt bU-\sqrt a\|_2^2=2(1-f),\qquad
 \|Z-a\|_1\le\sqrt{2(1-f)}\le\sqrt{2\eta}.
\]
Choose any rank-$r$ projection $P$ containing the range of $Z$.
This is possible since $\operatorname{rank}Z\le r$.  Since
$\tr Z=f$, one has
\[
 \|Z-fP/r\|_2^2=e,\qquad
 \|Z-fP/r\|_1\le\sqrt{re}.
\]
These formulas remain valid at $f=0$, without normalizing $Z$.
As $1-f\le\sqrt\eta$, the triangle inequality yields
\[
       \|a-P/r\|_1\le(1+\sqrt2)\sqrt\eta+\sqrt{re}.
\]
The normalized coefficient vectors for $\sqrt a$ and $\sqrt bU$
have inner product $f$ and receiver marginals $a,b$.
Their pure-state trace distance is $2\sqrt{1-f^2}$, by its
rank-two spectral calculation.  Contractivity under partial trace
therefore gives $\|a-b\|_1\le2\sqrt\eta$, the trace-distance/root-fidelity estimate of
Fuchs--van de Graaf~\cite[Theorem~1]{FvdG93}.  Consequently
\[
       \|b-P/r\|_1\le(3+\sqrt2)\sqrt\eta+\sqrt{re}.
\]
Weighted Cauchy--Schwarz bounds the square of either right side by
$[r+(3+\sqrt2)^2/c_{d,r}](c_{d,r}\eta+e)$.
Now use \eqref{eq:rankdefects93}.  Every argument used matrices on
finite supports and a unitary extension of a partial isometry; no
inverse of a singular matrix or limit of an equality statement occurs.
\end{proof}
The constants in this lemma are explicit but not asserted to be optimal.
The projection may vary between branches; the theorem does not force
all receivers to occupy one common subspace throughout the protocol.

\begin{corollary}[Near-optimal projection decompositions]
\label{cor:nearoptimalframes93}
Assume $(d,r)\ne(2,1)$ and a strategy in
$\mathfrak R_{k,\ell}$ has reward $v\ge p_r-\varepsilon$.
Choose a dimension-preserving refined normal form; on its nonzero
branches put
\[
 w_{yh}=\tr A_{yh},\qquad a_{yh}=A_{yh}/w_{yh},\qquad
 \mu_{yh}=w_{yh}/d,\qquad
 B_\varepsilon=\frac{2K_{d,r}d(d+1)}{t^2}\varepsilon.
\]
There are rank-$r$ projections $P_{yh}$ such that
\begin{align}
 \sum_{y,h}\mu_{yh}
  \bigl(\|a_{yh}-P_{yh}/r\|_1^2+
              \|b_{yh}-P_{yh}/r\|_1^2\bigr)&\le2B_\varepsilon,
 \label{eq:nearoptimalframes93}\\
 \frac1d\sum_y\left\|\rho-\sum_hw_{yh}P_{yh}/r\right\|_1^2
          &\le B_\varepsilon.
 \label{eq:framebarycenter93}
\end{align}
Finite and countable recorded instruments are allowed.  The weights
$\mu_{yh}$ sum to one and are Gram weights, not probabilities under a
chosen hypothesis.  For public mixtures the same statements hold with
the public seed adjoined to the refined record, using the common
purified initial Gram matrix.
\end{corollary}
\begin{proof}
Subtract \eqref{eq:unsaturatedrankupper93} from the attained class
value.  Since $\sum_{y,h}w_{yh}=d$, this gives
\[
 \sum_{y,h}\mu_{yh}\Delta_{d,r}(a_{yh},b_{yh})
       \le\frac{\varepsilon}{\alpha d}
       =\frac{2d(d+1)}{t^2}\varepsilon.
\]
Apply Lemma~\ref{lem:projectionstability93} to each branch and sum.
For the second estimate, each fixed $y$ has $\sum_hw_{yh}=1$ and
$\sum_hw_{yh}a_{yh}=\rho$.  The triangle inequality followed by
Jensen's inequality bounds its squared residual by
$\sum_hw_{yh}\|a_{yh}-P_{yh}/r\|_1^2$.
Average over $y$.  The proof uses nonnegative sums with finite total
weight, so countable traces pass to the limit.  The seed construction
in Lemma~\ref{lem:resetnormal91}, followed by the dimension-preserving
refinement, preserves the common Gram sum and the achieved reward.
\end{proof}
The total $\mu$-weight on branches for which either deviation is at
least $u>0$ is at most $2B_\varepsilon/u^2$.  This is a quantitative
statement at the relative loss scale $\varepsilon/t^2$, not a claim
of vanishing error at fixed nonzero calibration defect.

\subsection{A spectral loss certificate}
Let
\[
 \mathcal C_{d,r}=\{\sigma\in\mathcal D_d:\sigma\preceq I/r\},
 \qquad e_r(\rho)=\tr(\rho-I/r)_+.
\]
\begin{proposition}[Distance to the optimal initial spectra]
\label{prop:spectralloss93}
For every density matrix $\rho$,
\begin{equation}\label{eq:capdistance93}
 \inf_{\sigma\in\mathcal C_{d,r}}\|\rho-\sigma\|_1=2e_r(\rho).
\end{equation}
For $(d,r)\ne(2,1)$, any strategy in $\mathfrak R_{k,\ell}$
with initial Gram $\rho$ and achieved equal-prior reward $v$ satisfies
\begin{equation}\label{eq:spectralloss93}
 p_r-v\ge \frac{2t^2}{K_{d,r}d(d+1)}\,e_r(\rho)^2.
\end{equation}
\end{proposition}
\begin{proof}
Let $Q$ be the spectral projection of $\rho$ above $1/r$.
For $\sigma\preceq I/r$, one has
$\tr Q(\rho-\sigma)\ge e_r(\rho)$.
As $\tr(\rho-\sigma)=0$, positive-part optimization gives
$\|\rho-\sigma\|_1\ge2e_r(\rho)$.
Conversely, in an eigenbasis of $\rho$, cap every eigenvalue exceeding
$1/r$ and redistribute the removed mass among eigenvalues below $1/r$.
There is enough capacity since $d/r\ge1$.  This produces a density
matrix in $\mathcal C_{d,r}$ at trace distance exactly $2e_r(\rho)$.

Every matrix $\sum_hw_{yh}P_{yh}/r$ in
\eqref{eq:framebarycenter93} belongs to $\mathcal C_{d,r}$.
Set $\varepsilon=p_r-v\ge0$ there and use
\eqref{eq:capdistance93}; rearranging $4e_r(\rho)^2\le
2K_{d,r}d(d+1)(p_r-v)/t^2$ proves the bound.
\end{proof}
The estimate concerns the actual first-input marginal, up to its fixed
transpose, and a specified ideal reset architecture.  Under the
calibrated perturbations of Corollary~\ref{cor:gamestability91}, the
same strategy's nominal reward differs from its observed reward by at
most the stated channel, prior and latent-law error; a preparation
defect adds its stated half-trace contribution.  That total error
must be added to $p_r-v$ before applying the certificate.  There is no
uncalibrated inference of the reset cut from a score.

\subsection{The exceptional qubit face}
\begin{proposition}[Pinching at the exceptional endpoint]
\label{prop:qubitpinching93}
Let $d=2$, $r=1$, and suppose $a$ is pure.  Put
$\mathcal P_a(b)=aba+(I-a)b(I-a)$ and
$\Delta=\Delta_{2,1}(a,b)$.  Then
\begin{equation}\label{eq:qubitpinching93}
 \|b-\mathcal P_a(b)\|_1^2=2\Delta-\Delta^2\le2\Delta.
\end{equation}
In particular $\Delta=0$ if and only if $[a,b]=0$.
If instead $b$ is the pure factor, interchange $a,b$.
\end{proposition}
\begin{proof}
By simultaneous unitary conjugation take $a=|0\rangle\langle0|$ and
$b=\left(\begin{smallmatrix}x&z\\\bar z&1-x\end{smallmatrix}\right)$.
The nonzero block of the centered filtered swap is
$\sqrt b(I-2a)\sqrt b$.  Its trace is $1-2x$ and its determinant
is $-\det b$.  Its trace norm, also when $b$ is singular, is therefore
\[
 \sqrt{(1-2x)^2+4\det b}=\sqrt{1-4|z|^2}=1-\Delta.
\]
The pinching difference has eigenvalues $\pm|z|$.
Squaring gives \eqref{eq:qubitpinching93}.  Vanishing is exactly
$z=0$, or $[a,b]=0$.
\end{proof}
Thus the exceptional equality set includes a pure factor with any
commuting mixed qubit state, not just identical or orthogonal pure
factors.  A common rank-one projection cannot replace pinching here.
The initial-spectrum theorem remains valid because
$\mathcal C_{2,1}=\mathcal D_2$.  This exception is separate from the
zero-signal case $t=0$, in which every strategy has reward $1/2$.
